\documentclass{szb-solution}

\title{Elsőrendű és szétválasztható DE}
\area{Differenciálegyenletek}
\subject{Matematika G3}
\subjectCode{BMETE94BG03}
\date{Utoljára frissítve: \today}
\docno{9}

\usepackage{siunitx}
\sisetup{locale = DE, per-mode = symbol}

\usetikzlibrary{snakes}

\begin{document}
\maketitle

\clearpage

% ~~~~~~~~~~~~~~~~~~~~~~~~~~~~~~~~~~~~~~~~~~~~~~~~~~~~~~~~~~~~~~~~~~~~~~~~~~~~~~
% 333333333333333333333333333333333333333333333333333333333333333333333333333333
% ~~~~~~~~~~~~~~~~~~~~~~~~~~~~~~~~~~~~~~~~~~~~~~~~~~~~~~~~~~~~~~~~~~~~~~~~~~~~~~
\subsection{Szukcesszív approximáció I}

Számítsa ki az $y' = y$ differenciálegyenlet szukcesszív approximációjával
kapott első négy közelítő függvényt, ha a kezdeti feltétel $y(0) = 1$!

\begin{align*}
  y_0(x)        & = 1 \text,
  \\
  y_1(x)        & = y_0 + \int_{x_0}^x f(t; y_0) \dd t
  = 1 + \int_0^x 0 \dd t = 1 \text,
  \\
  y_2(x)        & = y_0 + \int_{x_0}^x f(t; y_1) \dd t
  = 1 + \int_0^x 1 \dd t = 1 + x \text,
  \\
  y_3(x)        & = y_0 + \int_{x_0}^x f(t; y_2) \dd t
  = 1 + \int_0^x (1 + t) \dd t = 1 + x + \frac{x^2}{2} \text,
  \\
  y_4(x)        & = y_0 + \int_{x_0}^x f(t; y_3) \dd t
  = 1 + \int_0^x \left( 1 + t + \frac{t^2}{2} \right) \dd t
  = 1 + x + \frac{x^2}{2} + \frac{x^3}{6} \text,
  \\
                & \vdots
  \\
  y_{\infty}(x) & = 1 + x + \frac{x^2}{2} + \frac{x^3}{6} + \cdots
  + \frac{x^k}{k!} + \cdots
  = \sum_{k=0}^{\infty} \frac{x^k}{k!} = e^x \text.
\end{align*}



% ~~~~~~~~~~~~~~~~~~~~~~~~~~~~~~~~~~~~~~~~~~~~~~~~~~~~~~~~~~~~~~~~~~~~~~~~~~~~~~
% 444444444444444444444444444444444444444444444444444444444444444444444444444444
% ~~~~~~~~~~~~~~~~~~~~~~~~~~~~~~~~~~~~~~~~~~~~~~~~~~~~~~~~~~~~~~~~~~~~~~~~~~~~~~
\subsection{Szukcesszív approximáció II}

Számítsa ki az $y' = xy$ differenciálegyenlet szukcesszív approximációjával
kapott első négy közelítő függvényt, ha a kezdeti feltétel $y(0) = 1$!

\begin{align*}
  y_0(x)        & = 1 \text,
  \\
  y_1(x)        & = y_0 + \int_{x_0}^x f(t; y_0) \dd t
  = 1 + \int_0^x t \cdot 1 \dd t = 1 + \frac{x^2}{2} \text,
  \\
  y_2(x)        & = y_0 + \int_{x_0}^x f(t; y_1) \dd t
  = 1 + \int_0^x t \left( 1 + \frac{t^2}{2} \right) \dd t
  = 1 + \frac{x^2}{2} + \frac{x^4}{8} \text,
  \\
  y_3(x)        & = y_0 + \int_{x_0}^x f(t; y_2) \dd t
  = 1 + \int_0^x t \left( 1 + \frac{t^2}{2} + \frac{t^4}{8} \right) \dd t
  = 1 + \frac{x^2}{2} + \frac{x^4}{8} + \frac{x^6}{48} \text,
  \\
                & \vdots
  \\
  y_{\infty}(x) & = 1 + \frac{x^2}{2} + \frac{x^4}{8} + \frac{x^6}{48} + \cdots
  + \frac{x^{2k}}{2^k k!} + \cdots
  = \sum_{k=0}^{\infty} \frac{x^{2k}}{2^k k!} = e^{x^2/2} \text.
\end{align*}



% ~~~~~~~~~~~~~~~~~~~~~~~~~~~~~~~~~~~~~~~~~~~~~~~~~~~~~~~~~~~~~~~~~~~~~~~~~~~~~~
% 666666666666666666666666666666666666666666666666666666666666666666666666666666
% ~~~~~~~~~~~~~~~~~~~~~~~~~~~~~~~~~~~~~~~~~~~~~~~~~~~~~~~~~~~~~~~~~~~~~~~~~~~~~~
\subsection{Szétválasztható differenciálegyenletek}

\begin{enumerate}[a)]
  \item $(2x + 1) y' - 3y = 0$
        $$
          \frac{\dd y}{y} = \frac{3 \dd x}{2x + 1}
          \quad \Rightarrow \quad
          \ln y = \frac{3}{2} \ln(2x + 1) + K
          \quad \Rightarrow \quad
          y = C (2x + 1)^{\sfrac{3}{2}}
        $$


  \item $\displaystyle\sqrt{1 + x^2} \, y' - \sqrt{1 - y^2} = 0$
        $$
          \frac{\dd y}{\sqrt{1 - y^2}} = \frac{\dd x}{\sqrt{1 + x^2}}
          \quad \Rightarrow \quad
          \arcsin y = \arcsinh x + K
          \quad \Rightarrow \quad
          y = \sin(\arcsinh x + K)
        $$

  \item $y' = \frac{1 - x - y}{2x - 2y - 3}$

        Írjuk át úgy az egyenletet, hogy $2x - 2y -3$ szerepeljen a nevezőben
        is:
        $$
          \frac{1 - x - y}{2x - 2y - 3}
          = -\frac12 \frac{2x -2y - 3 + 1}{2x - 2y - 3}
        $$

        Éljünk az $u = 2x - 2y - 3$ helyettesítéssel, ekkor:
        $$
          y' =
        $$

  \item $xy' = y(1 + \ln x - \ln y)$
  \item $2xyy' = x^2 + y^2$

  \item $y' = \frac{3x^2 + 4x + 2}{2y - 2}$, \quad $y(0) = -1$

        A differenciálegyenlet szétválaszható, így az alábbi megoldáshoz jutunk:
        $$
          (2y - 2) \dd y = (3x^2 + 4x + 2) \dd x
          \quad \Rightarrow \quad
          y^2 - 2y = x^3 + 2x^2 + 2x + C
          \text.
        $$
        A kezdeti feltétel alapján $y(0) = -1$, amiből kiszámítható a $C$ értéke:
        $$
          (-1)^2 - 2 \cdot (-1) = 0 + 0 + 0 + C
          \quad \Rightarrow \quad
          C = 3
          \text.
        $$
        Végül az implicit megoldás:
        $$
          y^2 - 2y = x^3 + 2x^2 + 2x + 3
          \text.
        $$

  \item $xy' = x \cdot e^{\sfrac{y}{x}} + y$, \quad $y(1) = 0$

        Éljünk az $u = y / x$ helyettesítéssel, ekkor
        $$
          y = x \, u
          \quad \text{és} \quad
          y' = u + x \, u'
          \text.
        $$
        Visszahelyettesítve az eredeti egyenletbe:
        $$
          x (u + x \, u') = x \, e^u + x \, u
          \quad \Rightarrow \quad
          x \, u' = e^u
        $$
        Az egyenlet szétválaszható, így az alábbi megoldáshoz jutunk:
        $$
          \frac{\dd u}{e^u} = \frac{\dd x}{x}
          \quad \Rightarrow \quad
          -e^{-u} = \ln x + C
          \text.
        $$
        A kezdeti feltétel alapján $y(1) = 0$, tehát $u(1) = 0$, amiből
        kiszámítható a $C$ értéke:
        $$
          -e^{0} = \ln 1 + C
          \quad \Rightarrow \quad
          C = -1
          \text.
        $$
        Végül visszahelyettesítve $u$-t:
        $$
          -e^{-\frac{y}{x}} = \ln x - 1
          \quad \Rightarrow \quad
          y = -x \ln(1 - \ln x)
          \text.
        $$
\end{enumerate}



% ~~~~~~~~~~~~~~~~~~~~~~~~~~~~~~~~~~~~~~~~~~~~~~~~~~~~~~~~~~~~~~~~~~~~~~~~~~~~~~
% 777777777777777777777777777777777777777777777777777777777777777777777777777777
% ~~~~~~~~~~~~~~~~~~~~~~~~~~~~~~~~~~~~~~~~~~~~~~~~~~~~~~~~~~~~~~~~~~~~~~~~~~~~~~
\subsection{Görbe egyenlete}

Adja meg azon görbét, amelynek bármely pontjában az érintő a
koordináta-ten\-ge\-lyek közé eső részét az adott pontban felezi!

Ilyen görbék esetén az $(x; y)$ pontban felírt érintő egyenes az $x$-tengelyt
a $(2x; 0)$ pontban, míg az $y$-tengelyt a $(0; 2y)$ pontban metszi. Az
egyenes meredeksége tehát:
$$
  m = \frac{2y - 0}{0 - 2x} = -\frac{y}{x}
  \text.
$$
Ez pont megegyezik $y'$-vel, így a következő differenciálegyenletet kapjuk:
$$
  y' = -\frac{y}{x}
  \text.
$$
A differenciálegyenlet szétválaszható, így az alábbi megoldáshoz jutunk:
$$
  \frac{\dd y}{y} = -\frac{\dd x}{x}
  \quad \Rightarrow \quad
  \ln |y| = -\ln |x| + K
  \quad \Rightarrow \quad
  y = \pm C \cdot \frac{1}{x}
  \text.
$$


% ~~~~~~~~~~~~~~~~~~~~~~~~~~~~~~~~~~~~~~~~~~~~~~~~~~~~~~~~~~~~~~~~~~~~~~~~~~~~~~
% 888888888888888888888888888888888888888888888888888888888888888888888888888888
% ~~~~~~~~~~~~~~~~~~~~~~~~~~~~~~~~~~~~~~~~~~~~~~~~~~~~~~~~~~~~~~~~~~~~~~~~~~~~~~
\subsection{Lehűlő kenyér}

Newton-törvénye értelmében ismert, hogy egy test hőmérsékletének változása a
környezet hőmérsékletével való különbséggel arányos. Egy kenyeret a $t = 0$
időpillanatban kiveszünk a $\SI{200}{\degreeCelsius}$-os sütőből, majd hűlni
hagyjuk. 20 perc után $\SI{60}{\degreeCelsius}$-ra hűl le. Mennyi idő múlva éri
el a kenyér hőmérséklete a $\SI{30}{\degreeCelsius}$-ot, ha a környezet
hőmérséklete $\SI{20}{\degreeCelsius}$?

Legyen a hőmérséklet $x(t)$, ekkor a Newton-törvény alapján:
$$
  \dot x = \alpha (x - x_k)
  \text,
$$
ahol $x_k = \SI{20}{\degreeCelsius}$ a környezet hőmérséklete, $\alpha < 0$
pedig az arányossági tényező. A differenciálegyenlet szétválaszható, így az alábbi
megoldáshoz jutunk:
$$
  \frac{\dd x}{x - x_k} = \alpha \dd t
  \quad \Rightarrow \quad
  x(t) = x_k + C \cdot e^{\alpha t}
  \text.
$$
Tudjuk, hogy $x(0) = 200$, tehát $C = 180$. Ezenkívül $x(20) = 60$, amiből
kiszámítható az $\alpha$ értéke:
$$
  60 = 20 + 180 \cdot e^{20\alpha}
  \quad \Rightarrow \quad
  \alpha = \frac{1}{20} \ln \frac{2}{9}
  \text.
$$
Végül megkeressük azt az időpontot, amikor a hőmérséklet eléri a
$\SI{30}{\degreeCelsius}$-ot:
$$
  30 = 20 + 180 \cdot e^{\frac{1}{20} \ln \frac{2}{9} \; t}
  \quad \Rightarrow \quad
  t = 20 \cdot \frac{\ln \frac{30 - 20}{180}}{\ln \frac{2}{9}} \approx \SI{38.43}{perc}
  \text.
$$



% ~~~~~~~~~~~~~~~~~~~~~~~~~~~~~~~~~~~~~~~~~~~~~~~~~~~~~~~~~~~~~~~~~~~~~~~~~~~~~~
% 999999999999999999999999999999999999999999999999999999999999999999999999999999
% ~~~~~~~~~~~~~~~~~~~~~~~~~~~~~~~~~~~~~~~~~~~~~~~~~~~~~~~~~~~~~~~~~~~~~~~~~~~~~~
\subsection{Keveredő sóoldat}

$\SI{100}{\kilogram}$ $\SI{10}{\percent}$-os sóoldatot tartalmazó edénybe
másodpercenként $\SI{10}{\liter}$ tiszta víz áramlik be. Mikor lesz a sóoldat
koncentrációja $\SI{5}{\percent}$-os, ha a keveredés azonnal megtörténik, és
ugyanilyen sebességgel folyik ki az edényből a keverék?

A be, illetve kiáramló térfogatáram $\dot V = \SI{10}{\liter\per\second}$, a
beáramló koncentráció $p_\text{be} = 0$, a kezdeti koncentráció pedig
$p_0 = \SI{0.1}{\kilogram\per\liter}$. A kívánt koncentráció
$p_k = \SI{0.05}{\kilogram\per\liter}$, vagyis $x = \SI{5}{\kilogram}$.
A be és kiáramló anyagmennyiség:
$$
  x_\text{be} = \dot{V} \cdot p_\text{be} = 0
  \quad \text{és} \quad
  x_\text{ki}(t)
  = \dot{V} \cdot p(t)
  = \dot{V} \cdot \frac{x(t)}{V_0}
  = 0,1 \cdot x(t)
$$
Az anyagmérleg alapján a következő differenciálegyenletet kapjuk:
$$
  \dot{x}(t) = x_\text{be} - x_\text{ki}(t)
  \quad \Rightarrow \quad
  \dot{x}(t) = -0,1 \cdot x(t)
  \text.
$$
A differenciálegyenlet szétválaszható, így az alábbi megoldáshoz jutunk:
$$
  \frac{\dd x}{x} = -0,1 \dd t
  \quad \Rightarrow \quad
  \ln x = -0,1 t + K
  \quad \Rightarrow \quad
  x(t) = C \cdot e^{-0,1 t}
  \text.
$$
A kezdeti feltétel alapján $x(0) = 10$, tehát $C = 10$. Végül megkeressük azt az
időpontot, amikor $x(t) = 5$:
$$
  5 = 10 \cdot e^{-0,1 t}
  \quad \Rightarrow \quad
  t = -10 \cdot \ln \frac{5}{10} \approx \SI{6.93}{\second}
  \text.
$$

\end{document}